% Ελληνική Περίληψη

\begin{abstract}\setlanguage{greek}
Η παρούσα διατριβή πραγματεύεται τον σχεδιασμό και την ανάπτυξη μιας συσκευής προγνωστικής συντήρησης, με στόχο την απόκτηση και ανάλυση δεδομένων. Το γνωστικό αντικείμενο της εργασίας εντάσσεται στη διάγνωση μηχανολογικών σφαλμάτων, στην ανάλυση δονήσεων, καθώς και στις τεχνολογίες του Internet of Things (IoT).

Η βασική πρόκληση που επιχειρεί να αντιμετωπίσει η εργασία είναι η δημιουργία ενός οικονομικού αλλά ταυτόχρονα αξιόπιστου συστήματος, το οποίο θα μπορεί να επεξεργάζεται δεδομένα επιτάχυνσης που προέρχονται από βιομηχανικές μηχανές. Ο τελικός στόχος είναι η πρόβλεψη πιθανών αστοχιών, όπως κακή ευθυγράμμιση, φθορά ρουλεμάν ή άλλα μηχανικά ελαττώματα, παρέχοντας στο προσωπικό συντήρησης χρήσιμες μετρήσεις, δείκτες και ιστορικά δεδομένα. Με τον τρόπο αυτό, μπορούν να λαμβάνονται τεκμηριωμένες αποφάσεις σχετικά με τις απαραίτητες ενέργειες συντήρησης.

Η μεθοδολογία βασίστηκε σε μια προσέγγιση μοντελοποίησης, με τη χρήση ενός τριαξονικού αναλογικού επιταχυνσιομέτρου (ADXL335) για την απόκτηση δεδομένων. Τα σήματα δόνησης καταγράφηκαν στο πεδίο του χρόνου και αναλύθηκαν με τον Γρήγορο Μετασχηματισμό Fourier (FFT), ώστε να εξαχθούν φάσματα συχνοτήτων. Τα φάσματα αυτά χρησιμοποιήθηκαν για τη σύγκριση της «υγιούς» λειτουργικής κατάστασης μιας μηχανής με πιθανές καταστάσεις βλάβης, με την ύπαρξη ορίου στο 1g να σηματοδοτεί ενδεχόμενο πρόβλημα.

Η υλοποίηση περιλάμβανε τον σχεδιασμό και την ολοκλήρωση ενός hardware setup με μικροελεγκτή Arduino για τη συλλογή δεδομένων από τον αισθητήρα και Raspberry Pi για την ανάλυση και επεξεργασία τους. Η επικοινωνία των δύο συσκευών πραγματοποιήθηκε μέσω του πρωτοκόλλου UART. Το ολοκληρωμένο σύστημα απέδειξε την ικανότητά του να συλλέγει δεδομένα δόνησης υψηλής ακρίβειας, να εκτελεί φασματική ανάλυση και να απεικονίζει τα αποτελέσματα με τρόπο κατανοητό.

Τα αποτελέσματα της μελέτης δείχνουν ότι είναι εφικτή η χρήση εξαρτημάτων χαμηλού κόστους και ευρείας διαθεσιμότητας για την ανάπτυξη ενός λειτουργικού εργαλείου προγνωστικής συντήρησης. Τα συμπεράσματα αναδεικνύουν τον καθοριστικό ρόλο της φασματικής ανάλυσης με FFT στη διάγνωση μηχανικών βλαβών, ενώ ταυτόχρονα υπογραμμίζουν τις δυνατότητες που προσφέρουν οι τεχνολογίες IoT στη βελτίωση στρατηγικών προληπτικής και προγνωστικής συντήρησης σε πλήθος βιομηχανικών εφαρμογών.

\vspace*{\fill}
\noindent{\large\bf{Λέξεις-κλειδιά:}}\\ 
Predictive Maintenance, Accelerometer, Vibration Analysis, Fast Fourier Transform (FFT), Internet of Things (IoT)
\end{abstract}







